\subsection{转置}

在本条中，假设 X 是纯 n 维的。

14.3.1. 设 \(\Omega = \det(\mathbf{T}(\mathbf{X})^{*})\)（参见 7.9.9）；这是一个在每点秩为 1、以 X 为底、\(\mathbf{C}^{r-1}\) 类的向量丛。我们有

\[
\Omega = \bigwedge^ {n} \mathrm{T} (\mathrm{X}) ^ {*} = \operatorname{Alt} ^ {n} (\mathrm{T} (\mathrm{X}); \mathrm{K} _ {\mathrm{X}}).
\]

设 M 是一个以 X 为底的向量丛。令

\[
\tilde {\mathbf {M}} = \mathcal {L} (\mathbf {M}; \Omega) = \mathbf {M} ^ {*} \otimes \Omega . ^ {(1)}
\]

若 M 是 \(C^{h}\) 类的，其中 \(h \in N_{K} \cup \{0\}\) 且 \(h \leqslant r - 1\)，则 \(\tilde{M}\) 也是。从 M 到 \(\tilde{\tilde{M}} = \mathcal{L}(\mathcal{L}(M; \Omega); \Omega)\) 的典范映射是一个同构；用它把 M 与 \(\tilde{M}\) 恒同。

14.3.2（“转置的定义”）。设 \(k\) 是一个 \(\leqslant r - 1\) 的正整数，并设 D 是 X 上的一个 E \(\rightarrow\) F 型、阶 \(\leqslant k\)、C \(^{h}\) 类的微分算子，其中 \(h \in \mathrm{N}_{\mathbf{K}} \cup \{0\}\) 且 \(k \leqslant h \leqslant r - k\)。那么存在一个微分算子 \(^{t}\) D（X 上的算子），它是 \(\widetilde{\mathbf{F}} \to \widetilde{\mathbf{E}}\) 型、阶 \(\leqslant k\) 的，具有下列性质：

设 \(\xi = (\xi^1, \ldots, \xi^n)\) 是 X 的开集 U 内的一个坐标系，并设 \(s = (s_i)_{1 \leqslant i \leqslant e}\)（相应地，\(t = (t_j)_{1 \leqslant j \leqslant f}\)）是 E（相应地，F）在 U 上的一个标架。设 \((s_i^*)\)（相应地，\((t_j^*))\)）是 E*（相应地，F*）的满足 \(\langle s_i, s_j^* \rangle = \delta_{ij}\)（相应地，\(\langle t_i, t_j^* \rangle = \delta_{ij}\)）的标架，并设 \((\tilde{s}_i)\)（相应地，\((t_j\)）是 \(\tilde{E}\)（相应地，\(\tilde{F}\)）在 U 上的标架，它由 \((s_i^*)\)（相应地，\((t_j^*)\)）与标架 \(\omega = d\xi^1 \wedge \cdots \wedge d\xi^n\)（\(\Omega\) 的标架）作张量积得到。设 \(c_\alpha^{ij}\)（\(1 \leqslant i \leqslant e\)，\(1 \leqslant j \leqslant f\)，\(|\alpha| \leqslant k\)）是 U 上的 \(C^h\) 类函数，使得有

\[
\mathrm{D} \left(\sum_ {i} f _ {i} s _ {i}\right) = \sum_ {i, j, \alpha} c _ {\alpha} ^ {i j} \Delta_ {\xi} ^ {\alpha} (f _ {i}). t _ {j}
\]

对任意函数 \(f_{i}\)（属于 \(\mathcal{C}^{h+k}(\mathbf{U})\)）成立（参见 14.1.4, (3')）。于是有

\[
{ } ^ { t } \mathbf { D } \left( \sum _ { j } g _ { j } \tilde { t } _ { j } \right) = \sum _ { i , j , \alpha } ( - 1 ) ^ { | \alpha | } \Delta _ { \xi } ^ { \alpha } ( c _ { \alpha } ^ { i j } g _ { j } ) \tilde { s } _ { i }
\]

对任意函数 \(g_{j}\)（属于 \(\mathcal{C}^{h+k}(\mathbf{U})\)）成立。

前面的性质（须对任意 \(\xi\)、\(s\) 与 \(t\) 验证）唯一地确定 \(^{t}\) D。\(^{(1)}\) 算子 \(^{t}\) D 称为 D 的转置；它是 \(C^{h-k}\) 类的。若 \(h \geqslant 2k\)，则 \(^{t}\) D 的转置有定义且等于 D（此时把 \(\tilde{E}\) 与 \(\tilde{F}\) 分别与 E 和 F 恒同）。

映射 \(\mathbf{D} \mapsto {}^t\mathbf{D}\) 是 K-线性的。若 \(k = 0\)，即若 \(\mathbf{D}\) 是 \(\mathbf{E}\) 到 \(\mathbf{F}\) 的态射，则 \({}^t\mathbf{D}\) 是态射 \(\mathcal{L}(\mathbf{D}; \mathrm{Id}_{\Omega})\)（从 \(\tilde{\mathbf{F}}\) 到 \(\tilde{\mathbf{E}}\) 的态射）。

\({}^{t}D\) 的符号借助同构

\[
\mathsf {T S} ^ {k} (\mathbf {T} (\mathbf {X})) \otimes \mathscr {L} (\mathbf {E}; \mathbf {F}) \rightarrow \mathsf {T S} ^ {k} (\mathbf {T} (\mathbf {X})) \otimes \mathscr {L} (\tilde {\mathbf {F}}; \tilde {\mathbf {E}})
\]

由 D 的符号得到，该同构是自同构 \((-1)^{k}\)（\(\mathsf{T}\mathsf{S}^{k}(\mathbf{T}(\mathbf{X}))\) 的自同构）与同构 \(u\mapsto\mathcal{L}(u;\mathrm{Id}_{\Omega})\)（从 \(\mathcal{L}(\mathbf{E};\mathbf{F})\) 到 \(\mathcal{L}(\tilde{\mathbf{F}};\tilde{\mathbf{E}})\) 上的同构）的张量积。

14.3.3（“复合的转置”）。记号与假设即 14.1.8 的那些，并设 \(D'\) 是 \(C^{h'}\) 类的、\(D''\) 是 \(C^{h''}\) 类的，其中 \(h',h'' \in \mathrm{N}_{\mathbf{K}} \cup \{0\}\)，\(h' \geqslant k' + 2k''\) 且 \(h'' \geqslant k'' + 2k'\)。

于是 \,、\, 与 \(\circ\) 的转置都有定义，\(^t\mathbf{D}'\) 与 \(^t\mathbf{D}''\) 的复合也有定义，并且我们有

\[
{ } ^ { t } ( \mathbf { D } ^ { \prime \prime } \circ \mathbf { D } ^ { \prime } ) = { } ^ { t } \mathbf { D } ^ { \prime } \circ { } ^ { t } \mathbf { D } ^ { \prime \prime } .
\]

特别地，在 14.3.2 的假设与记号下，若 \(f\) 是 X 上的 \(\mathbf{C}^{h'}\) 类函数且 \(2k \leqslant h'\)，则有 \({}^t(\mathbf{D} \circ f) = f \circ {}^t\mathbf{D}\)。

例。—— 取 X 为 \(K^{n}\) 的一个开子集；用 \(\xi^{1}, \ldots, \xi^{n}\) 表示 X 上的坐标函数，并把 \(\Omega\) 与 \(K_{x}\) 恒同（借助标架 \(\omega = d\xi^{1} \wedge \cdots \wedge d\xi^{n}\)）。若 \(E = F = K_{x}\)，则有

\[
\tilde {\mathbf {E}} = \tilde {\mathbf {F}} = \mathcal {L} (\mathrm{K}_{\mathbf{X}}; \Omega) = \mathcal {L} (\mathrm{K}_{\mathbf{X}}; \mathrm{K}_{\mathbf{X}}) = \mathrm{K}_{\mathbf{X}},
\]

且运算 \(\mathbf{D} \mapsto {}^{t}\mathbf{D}\) 把标量微分算子变为标量微分算子。例如有

\[
{ } ^ { t } ( \Delta ^ { \alpha } ) = ( - 1 ) ^ { | \alpha | } \Delta ^ { \alpha } \quad \text { for all } \alpha \in \mathbf {N} ^ { n } .
\]

若 \(r\) 是无穷，转置是代数 \(\mathcal{D}_{\mathbf{x}}^{\infty,r}(\mathrm{K}_{\mathbf{x}},\mathrm{K}_{\mathbf{x}})\) 的一个反自同构。

14.3.4. 设 \(k\) 是一个正整数，使得 K 的特征为 0 或 \(>k\)。设 \(\Omega_{1}\) 是向量丛 \(\bigwedge^{n-1}T(X)^{*} = \text{Alt}^{n-1}(T(X); K_{X})\)。记号与假设即 14.3.2 的那些，此外设 \(D'\) 是 X 上的一个 \(\widetilde{F} \to \widetilde{E}\) 型、阶 \(\leqslant k\) 的微分算子。我们称任何 \((E, \widetilde{F}) \to \Omega_{1}\) 型、阶 \(\leqslant k - 1^{(2)}\)、\(C^{h}\) 类（其中 \(h \geqslant 1\)）的多微分算子 G（参见 14.1.11），使得有恒等式

\[
\langle \mathrm{D} (u), v \rangle - \langle u, \mathrm{D} ^ {\prime} (v) \rangle = d (\mathrm{G} (u, v))\tag{1}
\]

对 X 的每个开集 U 及所有元素 \(u \in \mathcal{S}_{\mathrm{E}}^{k}(\mathrm{U})\)、\(v \in \mathcal{S}_{\mathrm{F}}^{k}(\mathrm{U})\) 成立。需要说明的是，在这个公式中，\(d\) 表示外微分（8.3.5），而 \(\langle D(u), v \rangle\)（相应地，

\(\langle u, \mathrm{D}'(v) \rangle\)）表示 \(\Omega\) 在 U 上的截面，它由对 \((\mathbf{D}(u), v)\)（相应地，对 \((u, \mathrm{D}'(v)))\)）经 \(F \times \tilde{F}\)（相应地，\(E \times \tilde{E}\)）到 \(\Omega\) 的典范配对得到。G 的存在性蕴含 \(D' = {}^t D\)；也称 G 为 D 的 Green 算子。

14.3.5. 设 D 是 X 上的一个 E \(\to\) F 型、阶 \(\leqslant k\)、\(\mathbf{C}^{h}\) 类的微分算子，其中 \(h \in N_{K} \cup \{0\}\) 且 \(k \leqslant h \leqslant r-k\)。在下列每种情形中，都存在 D 的一个 \(\mathbf{C}^{h-k+1}\) 类 Green 算子：

\begin{enumerate}
\def\labelenumi{\alph{enumi})}
\item
  \(\mathbf{K} = \mathbf{R}, r = \infty\) 且 X 是仿紧的。
\item
  K 的特征为 0 或 \(>k\)，X 同构于空间 \(K^{n}\) 的一个开子集，且丛 E 与 F 同构于平凡丛。
\item
  算子 D 的阶 \(\leqslant1\)（参见 14.3.7）。
\end{enumerate}

14.3.6. 我们保留 14.3.3 的假设与记号。若 G'（相应地，G''）是 D'（相应地，D''）的 Green 算子，则存在 D'' \(\circ\) D' 的 Green 算子 H，且是唯一的，使得有

\[
\mathbf {H} (u, v) = \mathbf {G} ^ {\prime \prime} (\mathbf {D} ^ {\prime} (u), v) + \mathbf {G} ^ {\prime} (u, ^ {t} \mathbf {D} ^ {\prime \prime} (v))
\]

对 X 的每个开集 U 及所有元素 \(u \in \mathcal{S}_{\mathrm{E}}^{k}(\mathrm{U})\) 与 \(v \in \mathcal{S}_{\widetilde{\mathrm{G}}}^{k}(\mathrm{U})\) 成立。

14.3.7. 我们有 \(\mathbf{S}^1 (\mathbf{E},\mathbf{F}) = \mathbf{T}(\mathbf{X})\otimes \mathbf{E}^*\otimes \mathbf{F}\)（参见 14.2.2）。另一方面，右内乘 \((\xi ,\omega)\mapsto i(\xi)\omega\)（参见 A, III, p. 158）定义了 \(\mathrm{T(X)}\otimes \Omega\) 到 \(\Omega_{1}\) 上的一个同构；与对偶 \(\Omega^{*}\)（\(\Omega\) 的对偶）作张量积，就得到 \(\mathrm{T(X)}\) 到 \(\Omega^{*}\otimes \Omega_{1}\) 上的一个同构，从而有向量丛的恒同

\[
\begin{array}{r l} \mathrm{S} ^ {1} (\mathrm{E}, \mathrm{F}) & = \Omega^ {*} \otimes \Omega_ {1} \otimes \mathrm{E} ^ {*} \otimes \mathrm{F} = (\mathrm{E} \otimes \mathrm{F} ^ {*} \otimes \Omega) ^ {*} \otimes \Omega_ {1} \\ & = (\mathrm{E} \otimes \tilde {\mathrm{F}}) ^ {*} \otimes \Omega_ {1} = \mathcal {L} (\mathrm{E} \otimes \tilde {\mathrm{F}}; \Omega_ {1}). \end{array}
\]

在此基础上，设 D 是 X 上的一个 E \(\to\) F 型、阶 \(\leqslant1\)、\(\mathbf{C}^{h}\) 类的微分算子，其中 \(h \in N_{K} \cup \{0\}\) 且 \(h \leqslant r-1\)。存在唯一的 D 的 Green 算子；它的阶为 0；它是 \(\mathbf{C}^{h}\) 类的、从 \(\mathbf{E}\otimes\widetilde{\mathbf{F}}\) 到 \(\Omega_{1}\) 的态射，并且由 D 的符号经上述恒同得到。

14.3.8. 假设 K = R，X 是分离且定向的（10.2.4），并设 A 是 X 的一个闭块（11.1.2）。给 A 装备由 X 的定向诱导的定向，并给 \(\partial A\) 装备相应的定向（11.2.1）。设 k 是满足 \(2k \leqslant r\) 的整数，D 是 X 上的一个 E \(\to\) F 型、阶 \(\leqslant k\)、\(C^{h}\) 类的微分算子。设 G 是 D 的一个 Green 算子。设 U 是 X 的开集，\(u \in \mathcal{S}_{\mathrm{E}}^{k}(\mathrm{U})\)，\(v \in \mathcal{S}_{\widetilde{\mathrm{F}}}^{k}(\mathrm{U})\)，并假设 u 与 v 的支集与 A 的交是紧集。于是有

\[
\int_ {\mathbf {A}} \langle \mathrm{D} (u), v \rangle - \int_ {\mathbf {A}} \langle u, ^ {t} \mathrm{D} (v) \rangle = \int_ {\partial \mathbf {A}} \mathrm{G} (u, v)
\]

（“Green 公式”）。特别地，若 \(\partial\mathbf{A}=\varnothing\)，则有

\[
\int_ {\mathrm{A}} \langle \mathrm{D} (u), v \rangle = \int_ {\mathrm{A}} \langle u, ^ {t} \mathrm{D} (v) \rangle .
\]


